\uuid{NgxZ}
\titre{Intégration par parties }

\niveau{} 				%L1, L2, L3, MPSI, MP, PCSI, PC, PSI...
\module{ Intégration } 	%Analyse, Algèbre...
\chapitre{Méthode d'intégration}   			%Continuité, Groupes, Fonctions de plusieurs variables...
\sousChapitre{Intégration par parties}			%Optimisation, Diagonalisation d'une matrice, Calcul de dérivées partielles...

\theme{}				%Fonction de répartition, Division euclidienne de polynômes, ...
\auteur{Erwan l'Haridon}
\datecreate{ 2026-02-11}
\organisation{AMSCC}			%AMSCC, Exo7, ...
\difficulte{2}			%1, 2, 3, 4 ou 5

\contenu{

\texte{ 

}

\begin{enumerate}
\item   \question{Calculer $\displaystyle \int_0^{1}\ln\!\left(x+\sqrt{x^2+1}\right)\,dx$.}
\indication{}
\reponse{On remarque $\ln\!\left(x+\sqrt{x^2+1}\right)=\operatorname{arsinh}(x)$. Par parties : $u=\operatorname{arsinh}(x)$, $dv=dx$, donc $du=\frac{dx}{\sqrt{x^2+1}}$, $v=x$. Ainsi $I=[x\,\operatorname{arsinh}(x)]_0^1-\int_0^1 \frac{x}{\sqrt{x^2+1}}dx$. Or $\int_0^1 \frac{x}{\sqrt{x^2+1}}dx=\sqrt{x^2+1}\big|_0^1=\sqrt2-1$. Donc $I=\ln(1+\sqrt2)-(\sqrt2-1)=\ln(1+\sqrt2)-\sqrt2+1$.}

\end{enumerate}

}